\section{The logic language}

To show the correspondence between semantics for lazy cost analysis and
functional logic programmming, we define a simple functional logic language $L$.
We show the syntax of $L$ in \cref{fig:syntax-logic}. \ly{TODO: Talk about all
  the new types we introduced and why we introduced them; talk about all the new
  terms as well.}

\begin{figure}[t]
\begin{center}
\begin{bnf}
\( A \) ::=
// \TBool
// \TProd{A}{A}
// \TThunk{A}
// \TNat
// \TList{A}
;;

$t$ ::=
| $x$
// \tlet{x}{t}{t}
// \ttrue{}
// \tfalse{}
// \teq{t}{t}
// \tle{t}{t}
// \tif{t}{t}{t}
| \tpair{t}{t}
// \tthunk{t}
// \tundefined{}
// \tcase{t}{x}{t}{t}
| $n$
// \tadd{t}{t}
// \tnil{}
// \tcons{t}{t}
// \tfoldr{x}{x}{t}{t}{t}
// \tchoose{t}{t}
// \tfree{A}
// \tfail{}
\end{bnf}
\end{center}
\caption{The syntax of the logic language $L$.}\label{fig:syntax-logic}
\end{figure}

\paragraph{Typing.}

\ly{Why is there only one rule? Are other rules mostly straightforward? Or you
  haven't got the time to write them down?}

\begin{center}
  \begin{mathpar}
    \inferrule{\TB{\Gamma, x_1 : A, x_2 : \TThunk{B}}{t_1}{B} \\
      \TB{\Gamma}{t_2}{B} \\
      \TB{\Gamma}{t_3}{\TList{A}}}
    {\TB{\Gamma}{\tfoldr{x_1}{x_2}{t_1}{t_2}{t_3}}{B}}
  \end{mathpar}
\end{center}

\paragraph{Semantics.}

\ly{TODO: Move a selected set of rules from \cref{apx:semantics-logic} here.}
\ly{Did we refer to any papers when we define this semantics? If so, cite them.}

\begin{lemma}
  If \( \TL{\Gamma}{t}{A} \), \( \gamma \in \Approx{\Gamma} \), and \( \EL{\gamma}{t}{v} \), then \( v \in \Approx{A} \).
\end{lemma}

Note that \( \tfail \) never reduces, so the language is not total.

\subsection{The clairvoyance translation}

We now define some constructs on the terms in \( \LL \).  For each construct \( c \), we provide typing rules and a set of evaluation rules that \emph{completely characterize} \( c \).  By ``completely characterize'', we mean the following.

\begin{itemize}
\item \emph{Soundness:} Each rule is valid.
\item \emph{Completeness:} Any derivation of the form \( \EL{\gamma}{c}{v} \) must have the form of one of the elimation rules for that construct.  In particular, there must be some elimination rule whose hypotheses are all true.  This allows us to ``invert'' judgments of the form \( \EL{\gamma}{c}{v} \).
\end{itemize}

\subsubsection{\texttt{return}}

\[
  \texttt{\treturn{t}} = \tpair{t}{0}
\]

\begin{center}
  \begin{prooftree}
    \hypo{\TL{\Gamma}{t}{A}}
    \infer1{\TL{\Gamma}{\treturn{t}}{\TM{A}}}
  \end{prooftree}
\end{center}

\begin{center}
  \begin{prooftree}
    \hypo{\EL{\gamma}{t}{v}}
    \infer1{\EL{\gamma}{\treturn{t}}{\spair{v}{0}}}
  \end{prooftree}
\end{center}

\subsubsection{bind}

\[
  \texttt{do $x$ <- $t_1$; $t_2$} = \tlet{\tpair{x}{y_1}}{t_1}{\tlet{\tpair{y_2}{y_3}}{t_2}{\tpair{y_2}{y_1 + y_3}}}
\]

\begin{center}
  \begin{prooftree}
    \hypo{\TL{\Gamma}{t_1}{\TM{A}}}
    \hypo{\TL{\Gamma, x : A}{t_2}{\TM{B}}}
    \infer2{\TL{\Gamma}{\texttt{do $x$ <- $t_1$; $t_2$}}{\TM{B}}}
  \end{prooftree}
\end{center}

\subsubsection{\texttt{transpose}}

\[
\begin{array}{l@{}l@{}l@{}}
  \texttt{transpose \(t\)} = \texttt{case \(t\) of } & \texttt{thunk \(y_1\) -> do $y_2$ <- $y_1$; return (\tthunk{y_2})} \\ & \texttt{undefined -> return undefined}
\end{array}
\]

\begin{center}
  \begin{prooftree}
    \hypo{\TL{\Gamma}{t}{\TThunk{(\TM{A})}}}
    \infer1{\TL{\Gamma}{\ttranspose{t}}{\TM{(\TThunk{A})}}}
  \end{prooftree}
\end{center}

\begin{center}
  \begin{prooftree}
    \hypo{\EL{\gamma}{t}{\sundefined}}
    \infer1{\EL{\gamma}{\ttranspose{t}}{\spair{\sundefined}{0}}}
  \end{prooftree}
\end{center}

\begin{center}
  \begin{prooftree}
    \hypo{\EL{\gamma}{t}{\sthunk{\spair{v}{c}}}}
    \infer1{\EL{\gamma}{\ttranspose{t}}{\spair{\sthunk{v}}{c}}}
  \end{prooftree}
\end{center}

\subsubsection{\texttt{foldM}}

\[
  \texttt{\tfoldM{x_1}{x_2}{t_1}{t_2}{t_3} = \tfoldr{x_1}{y_2}{\texttt{do \(x_2\) <- \ttranspose{y_2}; \(t_1\)}}{t_2}{t_3}} .
\]

\begin{center}
  \begin{prooftree}
    \hypo{\TL{\Gamma, x_1 : A, x_2 : \TThunk{B}}{t_1}{\TM{B}}}
    \hypo{\TL{\Gamma}{t_2}{\TM{B}}}
    \hypo{\TL{\Gamma}{t_3}{\TList{A}}}
    \infer3{\TL{\Gamma}{\texttt{foldM(\(\lambda x_1 x_2\).\(t_1\), \(t_2\), \(t_3\))}}{\TM{B}}}
  \end{prooftree}
\end{center}

To characterize the reduction behavior of \( \rulename{L-FoldM} \), we define yet another auxiliary relation:
\[
  \ELFM{\gamma}{x_1}{x_2}{t_1}{t_2}{v}{u} \equiv \ELF{\gamma}{x_1}{y_2}{\texttt{do \(x_2\) <- \ttranspose{y_2}; \(t_1\)}}{t_2}{v}{u} .
\]

\begin{lemma}
  The following rules completely characterize \( \tfoldMZ \).
  \begin{center}
    \begin{prooftree}
      \hypo{\ELFM{\gamma}{x_1}{x_2}{t_1}{t_2}{v}{u}}
      \infer1[\rulename{L-FoldM}]{\EL{\gamma}{\tfoldM{x_1}{x_2}{t_1}{t_2}{v}}{u}}
    \end{prooftree}
  \end{center}\begin{center}
    \begin{prooftree}
      \hypo{\EL{\gamma}{t_2}{u}}
      \infer1[\rulename{L-FoldM-Nil}]{\ELFM{\gamma}{x_1}{x_2}{t_1}{t_2}{\varepsilon}{u}}
    \end{prooftree}
  \end{center}
  \begin{center}
    \begin{prooftree}
      \hypo{\EL{\gamma, x_1 \mapsto v_1, x_2 \mapsto \sundefined}{t_1}{u}}
      \infer1[\rulename{L-FoldM-Undefined}]{\ELFM{\gamma}{x_1}{x_2}{t_1}{t_2}{\scons{v_1}{\sundefined}}{u}}
    \end{prooftree}
  \end{center}
  \begin{center}
    \begin{prooftree}
      \hypo{\ELFM{\gamma}{x_1}{x_2}{t_1}{t_2}{v_2}{\spair{v_3}{c_1}}}
      \hypo{\EL{\gamma, x_1 \mapsto v_1, x_2 \mapsto \sthunk{v_2}}{t_1}{\spair{v_4}{c_2}}}
      \infer2[\rulename{L-FoldM-Thunk}]{\ELFM{\gamma}{x_1}{x_2}{t_1}{t_2}{\scons{v_1}{\sthunk{v_2}}}{\spair{v_4}{c_1 + c_2}}}
    \end{prooftree}
  \end{center}
\end{lemma}
\begin{proof}
  We first prove validity.
  \newpage
  \newgeometry{margin=1cm}
  \begin{landscape}
    \begin{itemize}
    \item \textbf{\rulename{L-FoldM}.}
      \begin{center}
        \begin{prooftree}
          \hypo{\ELF{\gamma}{x_1}{y_2}{\texttt{do \(x_2\) <- \ttranspose{y_2}; \(t_1\)}}{t_2}{v}{u}}
          \infer1[\rulename{L-Foldr}]{\EL{\gamma}{\tfoldr{x_1}{y_2}{\texttt{do \(x_2\) <- \ttranspose{y_2}; \(t_1\)}}{t_2}{v}}{u}}
        \end{prooftree}
      \end{center}
    \item \textbf{\rulename{L-FoldM-Nil}.}
      \begin{center}
        \begin{prooftree}
          \hypo{\EL{\gamma}{t_2}{u}}
          \infer1[\rulename{L-Foldr-Nil}]{\ELF{\gamma}{x_1}{y_2}{\texttt{do \(x_2\) <- \ttranspose{y_2}; \(t_1\)}}{t_2}{\varepsilon}{u}}
        \end{prooftree}
      \end{center}
    \item \textbf{\rulename{L-FoldM-Undefined}.}
      \begin{center}
        \begin{prooftree}
          \infer0[\rulename{L-Var}]{\EL{\gamma, x_1 \mapsto v_1, y_2 \mapsto \bot}{y_2}{\bot}}
          \infer1[\rulename{L-Transpose-Undefined}]{\EL{\gamma, x_1 \mapsto v_1, y_2 \mapsto \bot}{\ttranspose{y_2}}{\spair{\bot}{0}}}
          \hypo{\EL{\gamma, x_1 \mapsto v_1, y_2 \mapsto \bot, x_2 \mapsto \sundefined}{t_1}{u}}
          \infer2[\rulename{L-Bind}]{\EL{\gamma, x_1 \mapsto v_1, y_2 \mapsto \bot}{\texttt{do \(x_2\) <- \ttranspose{y_2}; \(t_1\)}}{u}}
          \infer1[\rulename{L-Foldr-Undefined}]{\ELF{\gamma}{x_1}{y_2}{\texttt{do \(x_2\) <- \ttranspose{y_2}; \(t_1\)}}{t_2}{\scons{v_1}{\sundefined}}{u}}
        \end{prooftree}
      \end{center}
    \item \textbf{\rulename{L-FoldM-Thunk}.}  By induction, \( \ELF{\gamma}{x_1}{y_2}{\texttt{do \(x_2\) <- \ttranspose{y_2}; \(t_1\)}}{t_2}{v_2}{v_3} \).  Then,
      \begin{center}
        \begin{prooftree}
          \hypo{\ELF{\gamma}{x_1}{y_2}{\texttt{do \(x_2\) <- \ttranspose{y_2}; \(t_1\)}}{t_2}{v_2}{v_3}}
          \infer0[\rulename{L-Var}]{\EL{\gamma, x_1 \mapsto v_1, y_2 \mapsto \sthunk{\spair{v_3}{c_1}}}{y_2}{\sthunk{\spair{v_3}{c_1}}}}
          \infer1[\rulename{L-Transpose-Thunk}]{\EL{\gamma, x_1 \mapsto v_1, y_2 \mapsto \sthunk{\spair{v_3}{c_1}}}{\ttranspose{y_2}}{\spair{\sthunk{v_3}}{c_1}}}
          \hypo{\EL{\gamma, x_1 \mapsto v_1, y_2 \mapsto \sthunk{\spair{v_3}{c_2}}, x_2 \mapsto \sthunk{v_2}}{t_1}{\spair{v_4}{c_2}}}
          \infer2[\rulename{L-Bind}]{\EL{\gamma, x_1 \mapsto v_1, y_2 \mapsto \sthunk{\spair{v_3}{c_2}}}{\texttt{do \(x_2\) <- \ttranspose{y_2}; \(t_1\)}}{\spair{v_4}{c_1 + c_2}}}
          \infer2[\rulename{L-Foldr-Thunk}]{\ELF{\gamma}{x_1}{y_2}{\texttt{do \(x_2\) <- \ttranspose{y_2}; \(t_1\)}}{t_2}{\scons{v_1}{\sthunk{v_2}}}{\spair{v_4}{c_1 + c_2}}}
        \end{prooftree}
      \end{center}
    \end{itemize}
    To prove completeness, repeatedly apply inversion principles to find that any derivation of \( \tfoldM{x_1}{x_2}{t_1}{t_2}{v}{u} \) must have the form of one of the preceding derivations.
  \end{landscape}
  \restoregeometry
\end{proof}

\begin{center}
  \begin{align*}
    \floor{\tfoldr{x_1}{x_2}{t_1}{t_2}{t_3}} &= \texttt{do $y_3$ <- \(\floor{t_3}\); \tfoldM{x_1}{x_2}{\floor{t_1}}{\floor{t_2}}{y_3}}
  \end{align*}
\end{center}

\subsubsection{Properties of the translation}

\begin{lemma}
  If \( \TB{\Gamma}{t}{A} \), then \( \TL{\Gamma}{\floor{t}}{A} \).
\end{lemma}
